\todo[inline,color=gray!10]{meta: Dependencies: Section 02 (VCDM concepts, multi-level modeling, Refinery).}

\todo[inline,color=gray!10]{meta: Note: The teaser figure (fig\_teaser) is placed before the introduction in ACM sigconf. This section references it but does not carry the figure burden.}

\section{Motivation}\label{motivation}

\todo[inline,color=purple!15]{scaffold: M1 --- Housing subsidy scenario setup}

\todo[inline,color=purple!15]{scaffold: Job: Introduce the running example --- a real-world multi-issuer credential scenario that exercises all three layers.}

\todo[inline,color=purple!15]{scaffold: Key content: Government housing subsidy, 3 issuers, cross-credential dependency, privacy-sensitive income check, EU regulatory context.}

\todo[inline,color=purple!15]{scaffold: Key claim: Real credential ecosystems involve constraints from multiple governance sources spanning multiple abstraction layers.}

Consider a government housing subsidy where eligibility requires credentials from three independent authorities: family status from a civil registry, property records from a land registry, and income from an employer.\footnote{Based on the Hungarian Family Housing Subsidy (Családi Otthonteremtési Kedvezmény, CSOK), simplified. Additional credentials required in practice --- tax clearance, criminal record check --- are omitted.} The required property size depends on the number of children --- a constraint that spans two credentials --- and the income check must satisfy both EU format mandates and data protection requirements.

\todo[inline,color=blue!20]{cite: 518/2023. (XI. 30.) Korm. rendelet {\S}9 --- CSOK floor area requirements by number of children}

\todo[inline,color=purple!15]{scaffold: M2 --- Cross-layer constraints in the scenario}

\todo[inline,color=purple!15]{scaffold: Job: Walk through the scenario showing constraints at each layer and across layers.}

\todo[inline,color=purple!15]{scaffold: Key content: CPL has cross-property constraint (floor area depends on children). CSL has entity alignment (same Applicant across 3 credentials). FSL has governance conflict (eIDAS vs. GDPR on income credential).}

At the claim property layer, the applicant's facts --- number of children, property floor area, monthly income --- form an information graph with domain-level constraints: the minimum floor area is a function of the number of children \todo[color=blue!20]{cite: 518/2023 Korm. rendelet {\S}9}, and income must exceed a regulatory threshold. At the credential schema layer, these facts are distributed across three credentials issued by independent authorities, each with its own credential subject. All three subjects must be aligned --- they refer to the same applicant --- and each claim must trace to the corresponding domain-level property. At the format-specific layer, EU regulations mandate specific credential formats for government-issued attestations \todo[color=blue!20]{cite: eIDAS 2.0 ARF --- SD-JWT-VC/mdoc mandate}, while data protection law requires that privacy-sensitive checks --- such as whether income exceeds a threshold --- disclose only the minimum necessary information \todo[color=blue!20]{cite: GDPR Art. 5(1)(c), NAIH enforcement precedent}.

\todo[inline,color=purple!15]{scaffold: M3 --- Why single-layer inspection fails}

\todo[inline,color=purple!15]{scaffold: Job: Show that each layer is consistent in isolation but cross-layer constraints are violated.}

\todo[inline]{Draft --- single-layer inspection passes at each layer. Cross-layer analysis reveals two distinct problems: (1) the income credential's format cannot simultaneously satisfy eIDAS format mandates and GDPR data minimization --- SD-JWT-VC lacks predicate proofs, AnonCreds lacks VCDM conformance; (2) the floor area constraint requires cross-credential predicate evaluation, which no deployed format supports in zero-knowledge. Both problems are invisible when any single layer is inspected alone.}

\todo[inline,color=purple!15]{scaffold: M4 --- Governance framework complication}

\todo[inline,color=purple!15]{scaffold: Job: Show that constraints come from multiple, potentially conflicting governance sources.}

\todo[inline,color=blue!20]{cite: W3C VCDM 2.0 --- structural constraints (credentialSubject, proof types)}

\todo[inline,color=blue!20]{cite: eIDAS 2.0 ARF --- format mandate (SD-JWT-VC, mdoc)}

\todo[inline,color=blue!20]{cite: GDPR Art. 5(1)(c) --- data minimization}

\todo[inline,color=blue!20]{cite: AnonCreds specification --- CL-signature predicate proofs}

\todo[inline,color=blue!20]{cite: SD-JWT-VC --- IETF draft, hash-based selective disclosure}

\todo[inline]{Draft --- The constraints originate from governance frameworks that were not designed to be jointly satisfied: W3C VCDM defines structural conformance, eIDAS mandates credential formats, GDPR requires data minimization, and format-specific capabilities (predicate proofs, selective disclosure) vary across implementations. No existing tool or methodology checks whether these constraints can be simultaneously satisfied for a given credential ecosystem design.}

\todo[inline,color=purple!15]{scaffold: M5 --- Problem statement}

\todo[inline,color=purple!15]{scaffold: Job: State the problem that the rest of the paper addresses. End the section.}

\todo[inline,color=purple!15]{scaffold: Key claim: Credential ecosystem design requires a formal framework for simultaneous multi-layer, multi-source constraint satisfaction. No existing approach provides this.}

\todo[inline,color=purple!15]{scaffold: Note: This is where Sec 03 ends. The solution comes in Sec 04. Do NOT preview the solution here.}

\todo[inline,color=blue!20]{cite: gap claim --- no existing multi-level framework for VC ecosystems (confirmed by gap analysis 2026-03-25)}

\todo[inline]{Draft --- Credential ecosystem design requires a formal framework that captures constraints at multiple abstraction layers and from multiple governance sources, enabling designers to determine whether a given configuration is jointly satisfiable. No existing approach provides this capability.}
